\chapter{Future Work and Broader Reflections}
\label{chap:optional-discussion}

This chapter offers broader reflections on the findings of this study
and suggests directions for future research. While the core findings
and recommendations were presented in Chapters~3 and~4, this chapter
looks beyond the immediate scope of the thesis to consider wider
implications and open questions.

\section{Future Research Directions}
\label{sec:future}

This study provides a first empirical assessment of IS integration
maturity in PMO environments in Latvia. However, several important
research directions remain unexplored and could be pursued in
future work.

\textbf{Longitudinal research.} The current study captures a snapshot
of integration maturity at a single point in time. A longitudinal
study tracking the same organizations over several years would provide
valuable insights into how integration maturity evolves and what
factors drive or hinder progression between maturity levels. Such
research could also examine whether organizations that invest in
higher integration maturity sustain those gains over time, or whether
they regress when organizational priorities shift or key personnel
leave. Longitudinal data would further allow researchers to identify
the conditions under which organizations successfully transition from
Level 2 to Level 3 or Level 4, which the current cross-sectional
design cannot establish.

\textbf{Broader Baltic comparison.} Extending this study to include
Estonia and Lithuania would allow for a meaningful comparison of IS
integration maturity patterns across the Baltic region. Given that
Estonia is widely regarded as one of the most digitally advanced
countries in Europe, such a comparison could reveal important
differences and lessons for Latvian enterprises. Estonia's experience
with e-governance and digital infrastructure may have spillover
effects on enterprise-level IS integration that are not present in
Latvia or Lithuania, making a three-country comparison particularly
informative for policymakers and practitioners alike.

\textbf{Industry-specific benchmarks.} The current study treats all
industries as comparable. Future research could develop
industry-specific IS integration maturity benchmarks for Latvia,
recognizing that the integration needs and challenges of IT companies
differ significantly from those of manufacturing or public sector
organizations. For example, public sector PMOs face regulatory
constraints on data sharing and system procurement that private
sector organizations do not, which may fundamentally shape their
achievable integration maturity ceiling. Industry-specific benchmarks
would give organizations a more meaningful reference point against
which to evaluate their own maturity level.

\textbf{SME-focused research.} This study targeted enterprises with
50 or more employees. However, many Latvian enterprises are smaller
than this threshold. Future research could examine IS integration
maturity challenges specific to small and micro enterprises, which
face different resource constraints and may require tailored
integration solutions. For smaller organizations, the cost of
API-based integration or centralized BI platforms may be prohibitive,
making low-code and no-code solutions particularly relevant. Research
in this area could inform the development of affordable, scalable
integration pathways suited to the Latvian SME landscape.

\textbf{Role of emerging technologies.} The rapid development of
artificial intelligence, machine learning, and low-code platforms
is likely to significantly reshape IS integration in PMO environments
in the coming years. Future research could examine how these
technologies affect integration maturity progression and what new
maturity frameworks may be needed to capture their impact. In
particular, AI-driven data reconciliation and natural language
interfaces for BI dashboards may blur the boundaries between the
current IMM levels, potentially requiring a revised five- or
six-level model that accounts for intelligent automation as a
distinct capability tier.

\textbf{Validation of the Integration Maturity Model.} The four-level
IMM proposed and applied in this thesis was developed specifically
for the purposes of this study. While the empirical findings provide
initial validation of its discriminant properties, future research
should subject the model to more rigorous psychometric testing,
including confirmatory factor analysis and cross-national validation.
A validated, standardized IMM instrument would be a valuable
contribution to the IS research community and would enable
cumulative, comparable research across different organizational
and national contexts.

\section{Broader Societal and Economic Implications}
\label{sec:societal}

The findings of this study have implications that extend beyond
individual organizations. Latvia's digital transformation agenda,
supported by EU structural funds and national digitalization
strategies, places increasing emphasis on improving the digital
capabilities of Latvian enterprises. The finding that 70\% of PMO
environments operate at intermediate integration maturity levels
suggests that significant potential remains untapped.

Improving IS integration maturity in Latvian PMO environments could
generate wider economic benefits, including greater project delivery
efficiency, more effective use of public and private investment, and
stronger alignment between organizational strategy and project
execution. As Latvia continues to develop its digital economy,
investments in PMO IS integration may play an important role in
strengthening organizational competitiveness within the European
Union.

Furthermore, the recommendations proposed in Chapter~4 --- particularly
the emphasis on low-code platforms and data literacy training --- align
well with the EU's broader digital skills agenda and could be
incorporated into national programmes supporting the digital
transformation of Latvian enterprises. The European Commission's
Digital Decade targets, which aim for 75\% of EU enterprises to adopt
cloud, AI, or big data technologies by 2030, provide a broader policy
context within which the findings of this thesis are directly relevant.
Latvian enterprises that improve their PMO IS integration maturity are
not only strengthening their own operational effectiveness, but also
contributing to national progress toward these shared European digital
ambitions.

At the organizational level, the findings suggest that integration
maturity is not merely a technical concern but a governance and
strategic priority. PMO leaders and IT directors in Latvian enterprises
should treat IS integration maturity improvement as a multi-year
strategic initiative rather than a one-time technical upgrade. This
reframing --- from infrastructure project to strategic capability ---
is consistent with the broader shift in how digital transformation
is understood in contemporary management literature, and positions
integration maturity as a source of sustainable competitive advantage
rather than a compliance exercise.

\section{Chapter Summary}
\label{sec:chap5summary}

This chapter presented future research directions and broader
reflections on the implications of the study. Key directions for
future research include longitudinal studies, Baltic regional
comparisons, industry-specific benchmarking, SME-focused research,
investigation of the role of emerging technologies in PMO IS
integration, and formal psychometric validation of the Integration
Maturity Model. The findings of this study also carry wider societal
and economic implications for Latvia's digital transformation agenda,
connecting PMO IS integration maturity to the EU Digital Decade
targets and positioning integration improvement as a strategic
organizational priority rather than a purely technical undertaking.
\end{document}
